\chapter{Introducción y Objetivos}
Verificar de manera práctica el comportamiento de un Circuito Amplificador con Realimentación negativa.

\section{Conocimientos Previos}
\textbf{Unidad Temática 1:}\\
Realimentación negativa. Distintas topologías de circuitos realimentados. Cálculo de ganancia de lazo cerrado, lazo 
abierto, impedancias de lazo abierto y cerrado. Beneficios de circuitos realimentados. Concepto de desensibilidad.

\section{Equipamiento e Instrumental de Laboratorio}
\begin{itemize}
    \item Multímetro, Osciloscopio, Generador de Señales (Onda Senoidal)
    \item Fuente de alimentación variable hasta $+30 \V$
    \item Programa de simulación
    \item Hoja de Datos Transistor BC 337 o similar
\end{itemize}
